\notechapter{N-20220810}{Basic Inequalities, Cauchy, Triangle Inequality, and Coordinate Geometry}

\section{Basic inequalities}

The first page is a compact inequality note.  It begins with the elementary observation that squares are nonnegative.  For real numbers,
\[
  (a-b)^2\ge0,
\]
so
\[
  a^2+b^2\ge2ab,
\]
with equality iff $a=b$.  For three variables,
\[
  a^2+b^2+c^2\ge ab+bc+ca,
\]
because
\[
  2(a^2+b^2+c^2-ab-bc-ca)=(a-b)^2+(b-c)^2+(c-a)^2\ge0.
\]
The handwritten page marks equality conditions carefully: equality occurs when all variables involved in the squared differences coincide.

\section{AM--GM}

For $a,b\ge0$,
\[
  \frac{a+b}{2}\ge\sqrt{ab},
\]
with equality iff $a=b$.  Equivalently,
\[
  a+b\ge2\sqrt{ab}.
\]
The note also records the three-variable form
\[
  \frac{a+b+c}{3}\ge\sqrt[3]{abc},\qquad a,b,c\ge0,
\]
and uses it repeatedly in examples.

A typical application is to expressions such as
\[
  \frac{a+b+c}{3}\ge\sqrt[3]{abc}
  \quad\Longrightarrow\quad
  a+b+c\ge3\sqrt[3]{abc}.
\]
The useful habit is to look for terms whose product is fixed.  AM--GM then turns a fixed product into a lower bound for a sum.

\section{Cauchy's inequality}

The note states Cauchy's inequality in two forms.  The vector form is
\[
  (a_1b_1+\cdots+a_nb_n)^2
  \le
  (a_1^2+\cdots+a_n^2)(b_1^2+\cdots+b_n^2).
\]
The Engel or Titu Andreescu form is
\[
  \frac{x_1^2}{a_1}+\cdots+\frac{x_n^2}{a_n}
  \ge
  \frac{(x_1+\cdots+x_n)^2}{a_1+\cdots+a_n},\qquad a_i>0.
\]
This second form is often what the handwritten examples use.  It is obtained by applying Cauchy to
\[
  \left(\sum_i \frac{x_i^2}{a_i}\right)
  \left(\sum_i a_i\right)
  \ge
  \left(\sum_i x_i\right)^2.
\]

\section{Triangle inequality and its variants}

The page then discusses absolute values:
\[
  |a+b|\le |a|+|b|,\qquad
  |a-b|\le |a-c|+|c-b|.
\]
The second formula is the triangle inequality with an intermediate point $c$.  The note also uses the reverse triangle inequality
\[
  ||a|-|b||\le |a-b|.
\]
A useful proof is
\[
  |a|=|a-b+b|\le |a-b|+|b|,\qquad
  |b|=|b-a+a|\le |a-b|+|a|.
\]
Together these imply the reverse inequality.

The handwritten warning is that absolute value inequalities are not merely algebraic signs.  They encode distances on the real line.

\section{Worked AM--GM examples}

One example asks for a lower bound of expressions of the form
\[
  x+\frac1x,\qquad x>0.
\]
By AM--GM,
\[
  x+\frac1x\ge2.
\]
Similarly,
\[
  x+y+z\ge3\sqrt[3]{xyz}
\]
when $x,y,z>0$.  If a problem fixes $xyz$, this gives a direct lower bound for the sum.

Another common pattern is
\[
  \frac{x}{y}+\frac{y}{z}+\frac{z}{x}\ge3,\qquad x,y,z>0,
\]
because the product of the three terms is $1$.

\section{A Cauchy example}

For positive $a,b,c$, a typical inequality is
\[
  \frac{a^2}{b+c}+\frac{b^2}{c+a}+\frac{c^2}{a+b}
  \ge
  \frac{(a+b+c)^2}{2(a+b+c)}
  =\frac{a+b+c}{2}.
\]
This is exactly Engel form of Cauchy.  The handwritten page uses this kind of manipulation to turn a sum of fractions into one clean square divided by one clean denominator.

\section{Coordinate geometry at the end of the note}

The last page shifts to coordinate geometry.  The diagrams show points, vectors, and lines, with formulas for distances and slopes.  The essential formulas are:
\[
  AB=\sqrt{(x_A-x_B)^2+(y_A-y_B)^2},
\]
\[
  m=\frac{y_2-y_1}{x_2-x_1},
\]
and the point-slope form of a line
\[
  y-y_0=m(x-x_0).
\]

A recurring idea is that geometric claims can be turned into algebraic equalities.  For example, perpendicularity of two nonvertical lines is
\[
  m_1m_2=-1,
\]
while parallelism is
\[
  m_1=m_2.
\]
Distance and slope are the coordinate versions of length and direction.

\section{What the example reveals}

This note is about choosing the right inequality language.  Squares prove elementary inequalities.  AM--GM handles fixed products.  Cauchy handles sums of fractions and dot products.  Triangle inequalities handle distances.  Coordinate geometry then translates geometric distance and direction into algebra.  The common habit is to identify the hidden structure before calculating.

\section{Cauchy--Schwarz as projection}

Cauchy's inequality
\[
  \left(\sum_i a_i b_i\right)^2\le \left(\sum_i a_i^2\right)\left(\sum_i b_i^2\right)
\]
comes from inner-product geometry.  It says
\[
  |\langle a,b\rangle|\le \norm a\norm b.
\]
Equality holds exactly when the two vectors are linearly dependent.  Many contest equality cases are this geometric statement in coordinates.

\section{Triangle inequality and normed spaces}

The triangle inequality
\[
  \norm{x+y}\le \norm x+\norm y
\]
defines the geometry of a normed space.  For Euclidean norm it follows from Cauchy--Schwarz.  For \(L^p\) norms it becomes Minkowski's inequality.

Thus triangle inequality examples are early encounters with metric and norm geometry.

\section{AM--GM and convexity}

The inequality
\[
  \frac{x_1+\cdots+x_n}{n}\ge (x_1\cdots x_n)^{1/n}
\]
for positive variables follows from concavity of \(\log\):
\[
  \log\left(\frac1n\sum_i x_i\right)\ge \frac1n\sum_i\log x_i.
\]
The equality condition says all variables are equal, meaning the vector has no spread.

\section{Inequalities as geometry of norms and averages}

The basic inequalities are geometry and convexity in disguise.  Cauchy is projection, triangle inequality is norm structure, AM--GM is concavity of logarithm, and equality cases describe when the relevant geometric object degenerates to one direction or one value.

\section{Normed spaces and triangle inequality}

The triangle inequality in the original note is the defining feature of a norm.  A normed vector space is a vector space with a function $\|\cdot\|$ satisfying positivity, homogeneity, and
\[
  \|x+y\|\le \|x\|+\|y\|.
\]
The usual absolute value, Euclidean length, and sup norm are all instances.

Minkowski's inequality is the triangle inequality in $L^p$ spaces:
\[
  \|f+g\|_p\le \|f\|_p+\|g\|_p,\qquad
  \|f\|_p=\left(\int |f|^p\right)^{1/p}.
\]
Thus a contest-level inequality about square roots and sums is structurally the same statement as the triangle inequality for functions.

\section{Cauchy--Schwarz and Gram determinants}

Cauchy--Schwarz says
\[
  |\langle x,y\rangle|\le \|x\|\|y\|.
\]
In an inner product space, it follows from the nonnegativity of
\[
  \|x-ty\|^2\ge0.
\]
The equality case says that $x$ and $y$ are linearly dependent.  Geometrically, Cauchy--Schwarz bounds the length of a projection.

In Hilbert spaces, the same inequality controls Fourier coefficients:
\[
  |\langle f,e_n\rangle|\le \|f\|_2\|e_n\|_2.
\]
This is why an inequality first met in finite sums becomes indispensable in functional analysis.

\section{Convex cones and equality cases}

Many inequalities describe a cone of possible values.  For instance, the set of positive semidefinite matrices is a convex cone:
\[
  A\succeq0\quad\Longleftrightarrow\quad x^TAx\ge0\text{ for all }x.
\]
Cauchy--Schwarz can be read as positivity of the Gram matrix
\[
  \begin{pmatrix}
  \langle x,x\rangle&\langle x,y\rangle\\
  \langle y,x\rangle&\langle y,y\rangle
  \end{pmatrix}.
\]
Its determinant is nonnegative, giving the inequality.  Equality means the Gram matrix has rank one, i.e. the vectors are dependent.
